\subsection{Comparison table}

\begin{table}[ht]
\centering
\begin{tabularx}{\linewidth}{@{}>{\raggedright\arraybackslash}p{0.32\linewidth}
    >{\centering\arraybackslash}p{0.30\linewidth}
    >{\centering\arraybackslash}X@{}}
\toprule
Source & Value $(\times10^{-9})$ & Status \\
\midrule
Experimental discrepancy $\Delta a_\mu$ & $2.49\pm0.48$ & Fermilab + BNL world average \\
TFPT seam vertex $\delta_2/(2\pi)$ & $2.879$ & exact from carrier algebra \\
\midrule
Deviation from central value & $+0.389$ & $0.81\sigma$ \\
\bottomrule
\end{tabularx}
\caption{The TFPT prediction versus the measured muon $g-2$ discrepancy. The
predicted value lies within $0.81\sigma$ of the experimental central value and
well inside the $1\sigma$ band.}
\label{tab:comparison}
\end{table}

The TFPT prediction $a_\mu^{\mathrm{seam}}\approx\num{2.879e-9}$ overshoots
the central experimental value by
\[
\frac{a_\mu^{\mathrm{seam}}-\Delta a_\mu^{\mathrm{central}}}{\sigma_{\Delta a_\mu}}
=
\frac{(\num{2.879}-\num{2.49})\times10^{-9}}{\num{0.48e-9}}
\approx
0.81,
\]
placing it well inside the $1\sigma$ experimental band.

\subsection{Parameter-free comparison}

The comparison has zero adjustable parameters:

\begin{enumerate}[label=(\roman*), leftmargin=1.6em, itemsep=0.25em]
\item The carrier split $E_3\oplus E_2$ is the unique solution of
      $\dim\Lambda^{\mathrm{even}}E=16$ and $\tr_EY^2=5/6$.
\item The invariants $B=3/2$, $\gammaY=5/6$, and $B\gammaY=5/4$ follow
      algebraically.
\item The seam constant $\cthree=1/(8\pi)$ follows from the unit spectral
      flow normalization.
\item The occupancy $\Oadm=48$ is the topological family closure output.
\item The defect coefficients $\deltatop$ and $\delta_2$ are derived from
      (ii)--(iv).
\item The projection factor $1/(2\pi)$ is the inverse seam-loop phase.
\end{enumerate}

No parameter is fitted to the muon data.

\section{Why \texorpdfstring{$\delta_2$}{delta2} and not
\texorpdfstring{$\deltatop$}{delta-top}}\label{sec:why-delta2}

The first-order defect $\deltatop$ enters the exact seam generating function
$\phiseam(\alpha)$, which determines the fine-structure constant $\alpha$
through the carrier-form closure equation
\[
\alpha^3-2\cthree^3\alpha^2-8\cthree^6\bone\ln\!\bigl((\phiseam(\alpha))^{-1}\bigr)=0.
\]
By construction, $\deltatop$ is \emph{consumed} by the electromagnetic closure:
it is part of the input that fixes $\alpha$.

The second-order defect $\delta_2=\frac54\deltatop^2$ is the leading term in the
seam opening beyond what the electromagnetic closure absorbs. In the Taylor
expansion
\[
\phiseam(\alpha)
=
\phibase+\deltatop e^{-2\alpha}+\delta_2\,e^{-4\alpha}+O(e^{-6\alpha}),
\]
the $\deltatop$ term feeds the $\alpha$-equation, while $\delta_2$ constitutes
the first residual correction that manifests as a separate observable effect.

In standard QFT language, $\deltatop$ corresponds to the charge
renormalization, while $\delta_2/(2\pi)$ corresponds to the anomalous magnetic
moment---the leading vertex correction beyond the coupling itself.

\begin{proposition}[Order assignment]
\label{prop:order-assignment}
The electromagnetic closure absorbs the first-order defect $\deltatop$ into
the determination of $\alpha$. The anomalous magnetic moment is the
leading residual topological observable and is therefore governed by the
second-order defect $\delta_2$.
\end{proposition}


\section{Discussion}\label{sec:discussion}

\subsection{No new particles}

The TFPT derivation does not invoke supersymmetric partners (sleptons,
charginos, neutralinos), dark photons, leptoquarks, or extended Higgs
sectors. The anomaly is not a signal of new particles propagating in
loops; it is the topological signature of the seam---the codimension-one
locus where the involution-doubled manifold $\widetilde M$ folds onto
the physical quotient $M$.

\subsection{Universality}

The vertex projection $\delta_2/(2\pi)$ is lepton-universal at the topological
level: the seam defect does not depend on which lepton traverses the vertex.
This is consistent with the fact that the electron anomalous magnetic
moment $a_e$ is dominated by the QED Schwinger series and is measured to a
precision where the $\sim\!\num{2.9e-9}$ topological correction would appear
at the level of the $\alpha^5$ QED coefficient, which is currently at the
frontier of both theory and experiment.

\subsection{Lattice HVP and the residual tension}

Recent lattice QCD evaluations of the hadronic vacuum polarization
(HVP)~\cite{BMW2021,RBC2024} suggest a smaller discrepancy than the
dispersive approach. If the lattice central value is adopted, the
experimental anomaly shrinks to $\Delta a_\mu\approx(1.5\pm0.9)\times10^{-9}$
(current estimates). The TFPT prediction $\num{2.879e-9}$ would then sit
at $\sim\!1.5\sigma$, still within the allowed band but less tightly centered.
The resolution of the HVP tension between dispersive and lattice methods is
an experimental question; the TFPT prediction is fixed and testable against
either outcome.

\subsection{Falsifiability}

The prediction $a_\mu^{\mathrm{seam}}=45/(524288\,\pi^9)$ is exact and
non-negotiable within TFPT. If future measurements of $\Delta a_\mu$ converge
to a value outside the range $\num{2.879e-9}\pm\num{0.5e-9}$, the seam
vertex mechanism in its present form would be excluded.

\section{Conclusion}\label{sec:conclusion}

The muon $g-2$ anomaly admits a parameter-free explanation within the
TFPT carrier algebra. The second-order topological defect
\[
\delta_2
=
\frac54\,\deltatop^2
=
\frac{45}{262144\,\pi^8}
\approx\num{1.809e-8},
\]
projected through the seam-loop holonomy phase $\Delta_\Seam=2\pi$, yields
\[
\boxed{
a_\mu^{\mathrm{seam}}
=
\frac{\delta_2}{2\pi}
=
\frac{45}{524288\,\pi^9}
\approx\num{2.879e-9}.
}
\]
This lies within $0.81\sigma$ of the Fermilab/BNL world average discrepancy
$\Delta a_\mu=(2.49\pm0.48)\times10^{-9}$ and requires no supersymmetry, no
dark photons, and no free parameters.

The muon anomaly, in this reading, is not a messenger of new particles. It
is the topological signature of the seam---the intrinsic codimension-one
defect of the carrier-algebraic spacetime geometry.

\appendix

\section{Full derivation chain}\label{app:chain}

For self-containedness, we collect the complete algebraic chain from the
carrier split to the muon anomaly in one sequence.

\medskip

\textbf{Step 1: Carrier.}
$E=E_3\oplus E_2$ with $\dim\Lambda^{\mathrm{even}}E=16$,\quad
$\gammaY=\tr_EY^2=5/6$,\quad $B=3/2$.

\textbf{Step 2: Compression quotient.}
$B\gammaY=5/4$.

\textbf{Step 3: Seam constant.}
$\SF(U_\Seam)=1$, $\Delta_\Seam=2\pi$,
$\cthree=\Delta_\Seam/(16\pi^2)=1/(8\pi)$.

\textbf{Step 4: Occupancy.}
$N_{\mathrm{fam}}=3$, $\dim S^+=16$,
$\Oadm=N_{\mathrm{fam}}\dim S^+=48$.

\textbf{Step 5: First-order defect.}
$\deltatop=\Oadm\,\cthree^4=48/(8\pi)^4=3/(256\pi^4)\approx\num{1.203e-4}$.

\textbf{Step 6: Second-order defect.}
$\delta_2=B\gammaY\,\deltatop^2=\frac54\,\deltatop^2
=45/(262144\,\pi^8)\approx\num{1.809e-8}$.

\textbf{Step 7: Vertex projection.}
$a_\mu^{\mathrm{seam}}=\delta_2/\Delta_\Seam=\delta_2/(2\pi)
=45/(524288\,\pi^9)\approx\num{2.879e-9}$.

\textbf{Step 8: Comparison.}
$\Delta a_\mu^{\exp}=(2.49\pm0.48)\times10^{-9}$.
Deviation: $0.81\sigma$.


\section{Numerical verification}\label{app:numerics}

Independent numerical check using extended precision:

\begin{align}
\pi &= \num{3.14159265358979324}\notag\\
\pi^4 &= \num{97.4090910340024}\notag\\
256\pi^4 &= \num{24936.7273}\notag\\
\deltatop &= 3/\num{24936.7273} = \num{1.20304484e-4}\notag\\
\deltatop^2 &= \num{1.44731677e-8}\notag\\
\delta_2 &= \tfrac54\times\num{1.44731677e-8} = \num{1.80914597e-8}\notag\\
2\pi &= \num{6.28318531}\notag\\
a_\mu^{\mathrm{seam}} &= \num{1.80914597e-8}/\num{6.28318531}
= \num{2.8793e-9}\label{eq:check}
\end{align}


\section{Imported TFPT kernel block}\label{app:import}

The minimal import block for this application note is
\[
\bigl(
E_3\oplus E_2,\,
Y,\,
\inv,\,
\gammaY,\,
B,\,
\cthree,\,
\deltatop,\,
\delta_2,\,
\Oadm
\bigr).
\]
\begin{sloppypar}
This block is intentionally small enough to be verified
independently against the parent manuscript~\cite{TFPT31}
and the clean kernel note~\cite{TFPTkernel}.
\end{sloppypar}

\begin{sloppypar}
\begin{thebibliography}{99}

\bibitem{TFPT31}
S.~Hamann and A.~Rizzo,
\emph{TFPT: Carrier Algebra, Closure Theorems, and Intrinsic Observation},
Version~4.1 (2025).

\bibitem{TFPTkernel}
S.~Hamann and A.~Rizzo,
\emph{TFPT Minimal Kernel: Clean core note for downstream
application papers} (2025).

\bibitem{Fermilab2021}
B.~Abi \emph{et al.} (Muon $g-2$ Collaboration),
\emph{Measurement of the Positive Muon Anomalous Magnetic Moment to 0.46 ppm},
Phys.\ Rev.\ Lett.\ \textbf{126}, 141801 (2021).

\bibitem{Fermilab2023}
D.~P.~Aguillard \emph{et al.} (Muon $g-2$ Collaboration),
\emph{Measurement of the Positive Muon Anomalous Magnetic Moment to 0.20 ppm},
Phys.\ Rev.\ Lett.\ \textbf{131}, 161802 (2023).

\bibitem{BNL2006}
G.~W.~Bennett \emph{et al.} (Muon $g-2$ Collaboration),
\emph{Final Report of the Muon E821 Anomalous Magnetic Moment Measurement at BNL},
Phys.\ Rev.\ D~\textbf{73}, 072003 (2006).

\bibitem{WP2020}
T.~Aoyama \emph{et al.},
\emph{The anomalous magnetic moment of the muon in the Standard Model},
Phys.\ Rep.\ \textbf{887}, 1 (2020).

\bibitem{BMW2021}
Sz.~Borsanyi \emph{et al.} (BMW Collaboration),
\emph{Leading hadronic contribution to the muon magnetic moment from lattice QCD},
Nature \textbf{593}, 51 (2021).

\bibitem{RBC2024}
T.~Blum \emph{et al.} (RBC/UKQCD Collaboration),
\emph{Update of the hadronic vacuum polarization contribution to $a_\mu$},
Phys.\ Rev.\ Lett.\ \textbf{132}, 051903 (2024).

\end{thebibliography}
\end{sloppypar}

\end{document}
